% ch2_answers.tex -- Module 6, Chapter 2: answers to every Quick check question.
\begin{chapteranswers}{2}{The Motor Effect}
  % ---- Section: fields made by currents
  \item Change to metres: $r = \qty{4.0}{cm} = \qty{0.040}{m}$. Then
        \[ B = \frac{\mu_0 I}{2\pi r} = \frac{2 \times 10^{-7}\,\unit{N.A^{-2}} \times
        \qty{8.0}{A}}{\qty{0.040}{m}} = \qty{4.0e-5}{T}. \]
        Answer: \final{$B = \qty{4.0e-5}{T}$}.
  \item For a wire, $B \propto 1/r$. The new distance is $6.0 \div 2.0 = 3$ times bigger, so the
        field is 3 times smaller: $\qty{6.0e-5}{T} \div 3 = \qty{2.0e-5}{T}$.
        Answer: \final{$B = \qty{2.0e-5}{T}$}.
  \item Change to metres: $L = \qty{25}{cm} = \qty{0.25}{m}$. Then
        \[ B = \frac{\mu_0 NI}{L} = \frac{4\pi \times 10^{-7}\,\unit{N.A^{-2}} \times 500 \times
        \qty{1.2}{A}}{\qty{0.25}{m}} = \qty{3.02e-3}{T}. \]
        Answer: \final{$B = \qty{3.0e-3}{T}$} (2 significant figures).
  \item Curl the fingers of your right hand clockwise, the way the current goes. Your thumb points
        away from you, into the solenoid. So the field inside points away from you, and the field
        lines go \emph{into} the solenoid at this end. Field lines enter a magnet at its S end.
        Answer: \final{this end is S}. (A clockwise current seen from an end always means S.)
  % ---- Section: the motor effect
  \item (a) \final{\ang{90}}: the wire is perpendicular to the field, $\sin\ang{90} = 1$.
        (b) \final{\ang{0} or \ang{180}}: the wire is parallel to the field, $\sin\theta = 0$.
  \item $l = \qty{12}{cm} = \qty{0.12}{m}$ and $\theta = \ang{90}$, so
        $F = lIB = \qty{0.12}{m} \times \qty{4.0}{A} \times \qty{0.30}{T} = \qty{0.144}{N}$.
        Answer: \final{$F = \qty{0.14}{N}$}.
  \item $F = lIB\sin\theta = \qty{0.12}{m} \times \qty{4.0}{A} \times \qty{0.30}{T} \times
        \sin\ang{30} = \qty{0.144}{N} \times 0.5 = \qty{0.072}{N}$.
        Answer: \final{$F = \qty{7.2e-2}{N}$}, half of the force at \ang{90}.
  \item The electrons flow out of the page, so the conventional current flows \emph{into} the
        page. Palm rule: fingers to the left (field), thumb into the page (current). The palm
        faces up the page. Answer: \final{the force is up the page}.
  \item Rearrange $F = lIB$ (at \ang{90}) by dividing both sides by $lI$:
        \[ B = \frac{F}{lI} = \frac{\qty{4.0e-3}{N}}{\qty{0.050}{m} \times \qty{2.0}{A}}
        = \qty{0.040}{T}. \]
        Answer: \final{$B = \qty{0.040}{T}$} (\qty{40}{mT}).
  % ---- Section: parallel wires
  \item They \final{attract}. Picture both currents coming out of the page, with wire 1 on the
        left (\cref{ch2:fig:parallel}(a)). Grip rule: wire 1's field lines go anticlockwise, so at
        wire 2 the field points up the page. Wire 2 carries a current in this field, so it feels a
        motor-effect force. Palm rule: thumb out of the page, fingers up, palm faces left, towards
        wire 1. By Newton's third law, wire 1 feels an equal force towards wire 2. So the wires
        pull together.
  \item Change to metres: $r = \qty{2.0}{cm} = \qty{0.020}{m}$. Then
        \[ \frac{F}{l} = \frac{\mu_0}{2\pi}\frac{I_1 I_2}{r} = 2 \times 10^{-7}\,\unit{N.A^{-2}}
        \times \frac{\qty{5.0}{A} \times \qty{5.0}{A}}{\qty{0.020}{m}} = \qty{2.5e-4}{N.m^{-1}}. \]
        Answer: \final{$F/l = \qty{2.5e-4}{N.m^{-1}}$} (\qty{2.5e-4}{N} on each metre).
  \item $F/l$ is proportional to $I_1 I_2 / r$. Doubling one current doubles the top: $\times 2$.
        Halving $r$ halves the bottom, which doubles the force: $\times 2$. Total:
        $2 \times 2 = 4$. Answer: \final{$F/l$ is 4 times bigger}
        ($4 \times \qty{2.5e-4}{N.m^{-1}} = \qty{1.0e-3}{N.m^{-1}}$).
  \item By Newton's third law, the force of B on A and the force of A on B are equal in size.
        They are opposite in direction. The formula shows the same: both forces use the product
        $I_1 I_2 = 1.0 \times 4.0$, and the same $r$ and $l$. B's bigger current makes a bigger
        field at A, but A's current is smaller, so the two effects balance. Answer:
        \final{both wires feel the same size of force}.
  \item \emph{Definition:} take two wires. They are straight and parallel. They are infinitely
        long and very thin (negligible circular cross-section). They are \qty{1}{m} apart in a
        vacuum. The same steady current flows in each wire. If the force between the wires is $2 \times 10^{-7}$ N on each metre of length, the
        current in each wire is \qty{1}{A}.
        \emph{Link:} the definition names one force \emph{between} the wires. It does not say
        which wire. By Newton's third law, each wire pushes or pulls on the other with an equal and
        opposite force. So both wires feel exactly $2 \times 10^{-7}$ N per metre, and this one
        value can define the unit.
  % ---- Section: key experiments
  \item Change to kilograms: $\Delta m = \qty{0.35}{g} = \qty{3.5e-4}{kg}$. Then
        \[ F = \Delta m\, g = \qty{3.5e-4}{kg} \times \qty{9.8}{m.s^{-2}} = \qty{3.43e-3}{N}. \]
        Answer: \final{$F = \qty{3.4e-3}{N}$}.
  \item The rods are in series, so $I_1 = I_2 = \qty{10.0}{A}$. Change to metres:
        $r = \qty{5.0}{mm} = \qty{5.0e-3}{m}$. First find the force:
        \[ F = \frac{\mu_0}{2\pi}\frac{I_1 I_2\, l}{r} = 2 \times 10^{-7}\,\unit{N.A^{-2}} \times
        \frac{\qty{10.0}{A} \times \qty{10.0}{A} \times \qty{0.50}{m}}{\qty{5.0e-3}{m}}
        = \qty{2.0e-3}{N}. \]
        Then change the force to a balance reading:
        \[ \Delta m = \frac{F}{g} = \frac{\qty{2.0e-3}{N}}{\qty{9.8}{m.s^{-2}}}
        = \qty{2.04e-4}{kg} = \qty{0.204}{g}. \]
        Answer: \final{$\Delta m = \qty{0.20}{g}$}. The reading rises, because opposite currents
        repel and the lower rod is pushed down.
  \item The rods are in series, so both carry the same current $I$, and $I_1 I_2 = I^2$. Then
        \[ F = \frac{\mu_0}{2\pi}\frac{I^2 l}{r} = \left(\frac{\mu_0 l}{2\pi r}\right) I^2. \]
        $\mu_0$, $l$ and $r$ stay the same, so $F \propto I^2$: a straight line through the origin.
        Answer: \final{gradient $= \mu_0 l/(2\pi r)$} (here $\qty{2.0e-5}{N.A^{-2}}$ in theory).
  \item Same-direction currents attract. The upper rod is fixed, so the lower rod is pulled
        \emph{up}, towards it. It presses less on the balance. The currents are the same as in
        \cref{ch2:qc:rods}, so the force has the same size, \qty{2.0e-3}{N}. Answer: \final{the
        reading falls below zero}, to about $-\qty{0.20}{g}$.
\end{chapteranswers}
